\section{Promienie równoległe graniczne i końce} % lang-pl
\label{sekcja_promienie_graniczne}

Równoległość graniczna była dla Łobaczewskiego, Bolyaia i Gaussa podstawowym narzędziem, ale ani oni, ani Eves nie dowodzą wszystkiego, co potrzeba. % lang-pl
W tej sekcji zobaczymy, że równoległość graniczna jest relacją równoważności, i jak z jej klas zrodzi się pojęcie końca prostej -- punktu idealnego. % lang-pl

Hartshorne oznacza półprostą symbolem $Aa$, gdzie $A$ jest jej początkiem, a $a$ oznacza prostą ją zawierającą wraz z wyborem kierunku; dwie półproste są \emph{współkońcowe} (ang.~\emph{coterminal}), jeśli leżą na jednej prostej i „idą w tę samą stronę'' \cite[s. 312]{hartshorne_2000}. % lang-pl

\begin{definition}
[półproste równoległe graniczne] % lang-pl
\index{półproste!równoległe graniczne} % lang-pl
	Półprosta $Aa$ jest równoległa graniczna do półprostej $Bb$, jeśli są współkońcowe albo leżą na różnych prostych różnych od $AB$, nie przecinają się i każda półprosta wewnątrz kąta $BAa$ przecina półprostą $Bb$. % lang-pl
\end{definition}

Hartshorne \cite[s. 312]{hartshorne_2000} (notacja $Aa \parallel\!\!\parallel Bb$). % lang-pl

Istnienie takich półprostych nie wynika z aksjomatów neutralnych; zakłada je dopiero aksjomat Hilberta z sekcji \ref{sekcja_hiperboliczna} \cite[s. 313]{hartshorne_2000}. % lang-pl

\begin{proposition}
	Równoległość graniczna półprostych jest przenoszalna na półproste współkońcowe, symetryczna i przechodnia. % lang-pl
\end{proposition}

Dowody: Hartshorne \cite[s. 313--316]{hartshorne_2000} (prop. 34.9--34.11 i lemat 34.12); u Evesa te same fakty są w sekcji \ref{sekcja_hiperboliczna}. % lang-pl

\begin{definition}
[koniec] % lang-pl
\index{koniec prostej} % lang-pl
	Klasę równoważności półprostych równoległych granicznych nazywamy \emph{końcem}. % lang-pl
\end{definition}

Hartshorne \cite[s. 316]{hartshorne_2000} (wniosek 34.13). % lang-pl

Końce to dokładnie punkty idealne z sekcji \ref{sekcja_punkty_idealne}; Hilbert wyprowadzi z nich całe ciało liczbowe, a modele Kleina i Poincarego okażą się jedynymi płaszczyznami hiperbolicznymi (zobacz sekcję \ref{sekcja_ciala_koncow}). % lang-pl

Zadania: konstrukcje trójkąta granicznego i jego „dwusiecznych'' -- Hartshorne \cite[s. 317]{hartshorne_2000} (zad. 34.9--34.12). % lang-pl
% Źródła: Hartshorne s. 373--384 (§40); Greenberg 2010, s. 209, 211.

